\begin{itemize}
    \item Generate reference data for all incident angles \\ This is done by rotating the grid overlay on simulation data and then performing the binning described below
    \item Bin prompt and delayed displacement vectors in matrix form by counting the number of voxels separating the IBD vertex and the capture vertex. Normalize all matrices for reference angles
    \item Create binning matrix for data set and normalize it
    \item Compute the \(L^2\) norm for the matrix differences for all reference angles
    \item Fit \(\alpha\left|\sin{\left(\frac{\theta-\theta_0}{2}\right)}\right|\) to \(L^2\) norms as a function of angle. The minimum gives the \(\theta_0\) fit parameter for the direction of the neutrino source.
\end{itemize}

\vspace{2cm}

\begin{figure}[ht]
        \setlength{\unitlength}{.04\linewidth}
        
    \begin{minipage}[t]{0.48\linewidth}
        \centering
\begin{picture}(10,5)
  \footnotesize
        \multiput(5,0)(.5,0){9}{\color{black}\line(0,1){4}} 
        \multiput(5,0)(0,0.5){9}{\color{black}\line(1,0){4}} 

        \multiput(5,4)(.5,0){9}{\color{black}\line(1,1){1}} 
        \multiput(9,0)(0,.5){9}{\color{black}\line(1,1){1}} 

        \put(10,1){\line(0,1){4}}
        \put(6,5){\line(1,0){4}}

        
          \thicklines
        \multiput(6.5, 2.50)(.025,0){20}{\color{black}\line(0,1){0.5}} 
        \multiput(0, 5)(.025,0){20}{\color{black}\line(0,-1){0.5}} 
        
        \multiput(6, 1.50)(.025,0){20}{\color{gray}\line(0,1){0.5}} 
        \multiput(0, 4.25)(.025,0){20}{\color{gray}\line(0,-1){0.5}} 
          \thinlines

        \put(0.8,4.7){prompt segment}
        \put(0.8,3.9){delayed segment}


        \put(1,1){\vector(-1,-1){1}} 
        \put(1,1){\vector(0,1){2}} 
        \put(1,1){\vector(1,0){2}} 

        \put(1,1){\vector(2,1){2}} 
        \put(1,1){\arc[0,26]{1}} 
        \put(1,1){\arc[0,26]{.95}} 

        
        \put(0,0.3){$z$}
        \put(0.7,2.8){$y$}
        \put(2.8,.7){$x$}
        \put(2.7,2){$\nu$}
        
       \put(6.75,-0.1){\color{black!90}\vector(0,1){4.7}} 
        \multiput(6.7,0.)(0,.25){16}{\color{blue}\line(1,0){0.1}}
        \put(4.8,2.75){\color{black!90}\vector(1,0){4.8}} 
        \multiput(5,2.7)(.25,0){16}{\color{blue}\line(0,1){0.1}} 

        \put(6.8,4.55){\color{black}$y'$}
        \put(9.2,2.4){\color{black}$x'$}
        \put(5.85,2.1){\color{black}$\theta_n^i$}


        \put(6.75,2.75){\color{blue}\arc[0,245]{0.65}} 
        \put(6.75,2.75){\color{blue}\vector(-1,-2){0.5}} 

        
        \multiput(3.5, 0.5)(0,1){4}{\color{black}\vector(2,1){1}} 
        
       \put(2,0){neutrino wind} 
        \put(2.2,1.2){$\vartheta_\nu$}
\end{picture}

    \end{minipage}
    \hfill
    \begin{minipage}[t]{0.48\linewidth}
        \centering

%        \setlength{\unitlength}{.1\linewidth}
\begin{picture}(10,4)
  \footnotesize

        \multiput(3,0.5)(.5,0){6}{\color{black}\line(0,1){2.5}} 
        \multiput(3,0.5)(0,0.5){6}{\color{black}\line(1,0){2.5}} 


        
          \thicklines
        \multiput(4., 1.50)(.025,0){20}{\color{black}\line(0,1){0.5}} % vertical
          \thinlines

 
        \put(4.25,0){\color{black}\vector(0,1){4}} 
        \put(2.5,1.75){\color{black}\vector(1,0){5}} 
        \multiput(4.2,0.25)(0,.5){7}{\color{blue}\line(1,0){0.1}}
        \multiput(2.75,1.7)(.5,0){7}{\color{blue}\line(0,1){0.1}} 

        \put(3.9,3.7){\color{black}$y$}
        \put(7.3,1.4){\color{black}$x$}


        \put(4.25,1.75){\color{blue}\arc[0,45]{2}} 
        \put(2.5,0){\color{blue}\vector(1,1){4}} 
        \put(5.7,3.7){\color{black}$45^\circ$}

        \put(2.5,0.5){\color{blue}\vector(7,5){5}} 
        \put(2.5,1){\color{blue}\vector(7,3){5}} 
        \put(2.5,1.5){\color{blue}\vector(7,1){5}} 

        \put(2.5,1.25){\color{red}\vector(7,2){5}} 
        \put(7.5,2.5){\color{red}$\theta_\nu^{exp}$}
         \put(7.5,2.){\color{blue}$\theta_\nu^{MC_1}$}
        \put(7.5,3){\color{blue}$\theta_\nu^{MC_2}$}
        \put(7.5,3.7){\color{blue}$\theta_\nu^{MC_3}$}
     

\end{picture}
\end{minipage}
        \caption{Diagram explaining the prompt-delayed segment geometry, as well as neutrino and axes orientation. 
%        In this work we assume that neutrino ``wind'' is in $xy$ plane, at an angle $\nu$ with the $x$ axis.
%        The prompt and delayed segment in this diagram shown for illustrative purposes; the actual separation depends on the segment size --- for large segments, prompt and delayed events would happen in the same segment. It is also illustrated that capture is not necessarily aligned with the neutrino beam (only on average it is the case as discussed in this paper).
%        To study the angular dependence we vary the angle of neutrino with respect to the detector's $x$ and $y$ axes (i.e. rotation of the detector along $z$ axis).}
}
        \label{fig:phi_def}
\end{figure}















% \begin{figure}[ht!]
%         \setlength{\unitlength}{.1\linewidth}
% \begin{picture}(10,4)

%         \multiput(3,0.5)(.5,0){6}{\color{black}\line(0,1){2.5}} 
%         \multiput(3,0.5)(0,0.5){6}{\color{black}\line(1,0){2.5}} 


        
%           \thicklines
%         \multiput(4., 1.50)(.025,0){20}{\color{black}\line(0,1){0.5}} % vertical
%           \thinlines

 
%         \put(4.25,0){\color{black}\vector(0,1){4}} 
%         \put(2.5,1.75){\color{black}\vector(1,0){5}} 
%         \multiput(4.2,0.25)(0,.5){7}{\color{blue}\line(1,0){0.1}}
%         \multiput(2.75,1.7)(.5,0){7}{\color{blue}\line(0,1){0.1}} 

%         \put(3.9,3.7){\color{black}$y$}
%         \put(7.3,1.4){\color{black}$x$}


%         \put(4.25,1.75){\color{blue}\arc[0,45]{2}} 
%         \put(2.5,0){\color{blue}\vector(1,1){4}} 
%         \put(5.7,3.7){\color{black}$45^\circ$}

%         \put(2.5,0.5){\color{blue}\vector(7,5){5}} 
%         \put(2.5,1){\color{blue}\vector(7,3){5}} 
%         \put(2.5,1.5){\color{blue}\vector(7,1){5}} 

%         \put(2.5,1.25){\color{red}\vector(7,2){5}} 
%         \put(7.5,2.5){\color{red}$\theta_\nu^{exp}$}
%          \put(7.5,2.){\color{blue}$\theta_\nu^{MC_1}$}
%         \put(7.5,3){\color{blue}$\theta_\nu^{MC_2}$}
%         \put(7.5,3.7){\color{blue}$\theta_\nu^{MC_3}$}
     

% \end{picture}
% \caption{A diagram explaining the idea behind a pattern-matching algorithm for IBD segmented detectors. The experimental/empirical data set (with {\it a priori} unknown neutrino angle $\theta_\nu^{exp}$) is compared against theoretical/simulation data sets (of known neutrino angles $\theta_\nu^{MC_i}$). Due to the symmetry and square geometry unique distributions only obtained between 0$^\circ$ and 45$^\circ$ angles. MC$_1$, MC$_2$, and MC$_3$ denote Monte Carlo data sets corresponding to that particular neutrino angle. The 5$\times$5 grid is shown for illustrative purposes.}
% \label{fig_algo_skewers}
% \end{figure}